% =============================================================================
% Chapter 7: Discussion
% Interprets the findings of Chapter 6. Items marked as open checks in the
% handover are presented as interpretations, not as findings.
% =============================================================================

\chapter{Discussion}
\label{ch:discussion}

Chapter~\ref{ch:results} states what the register does on 94,816 simulated left
turns. This chapter asks what those findings establish and what they do not.

Section~\ref{sec:disc_scope} fixes the scope of the claim. The three sections
that follow interpret the main findings: why leading indicators are needed at all
(Section~\ref{sec:disc_leading}), what separates an indicator that tracks the
configuration from one that tracks exposure (Section~\ref{sec:disc_exposure}),
and why the conflict geometry decides the warning time
(Section~\ref{sec:disc_geometry}). Sections~\ref{sec:disc_pedrule}
to~\ref{sec:disc_threshold} discuss four individual decisions and anomalies: the
pedestrian collision rule, the behaviour of AEB\_PED and the reaction tier, the
formulation of REQ\_DECEL, and the choice of metric threshold.
Section~\ref{sec:disc_limitations} collects the limitations.

% =============================================================================
\section{What the Results Validate}
\label{sec:disc_scope}

The scope of the claim has to be stated before the findings are interpreted,
because the evaluation is easily read as saying more than it does.

The results validate the derivation method and the evaluation, not the safety of
a driving function. The ego vehicle is driven by the CARLA Traffic Manager with a
pedestrian-stop rule and without gap acceptance, so the conflicts that arise in
the oncoming chain follow from a capability the vehicle does not have rather than
from a functional insufficiency of a system that has it. The indicators are
computed from logged ground truth rather than from a perceived world model. What
the evaluation therefore shows is that the signals the method derives behave as
precursors of the collisions they were derived from, under a controlled and
reproducible exposure. It shows nothing about how safe the simulated vehicle is,
and a collision count obtained under a generator built to provoke collisions
carries no information about collision frequency in operation.

% =============================================================================
\section{Why Leading Indicators Are Needed}
\label{sec:disc_leading}

The argument for leading indicators is usually made in principle, by observing
that losses are too rare to count. The two configurations evaluated here allow it
to be made with measurements instead, because they place the same register in
both regimes.

In Campaign~1 the leading SPIs and the collisions are of comparable size. The
oncoming chain records 3.8 GAP\_TTC violations, 1.9 REQ\_DECEL violations and 1.2
AEB\_ONC events per collision, and the pedestrian chain 1.6 PED\_CPA violations
and 1.0 AEB\_PED events per collision. In that regime the lagging indicator is
itself countable, and a monitoring process would gain little from the leading
one.

In Campaigns~2 and~3 the position reverses. Across 75,509 runs a single collision
occurs, while the leading SPIs continue to be violated between 527 and 9,841
times. A process that monitored only the lagging indicator would have almost
nothing to observe, whereas every leading SPI still produces a measurable
quantity. This is the regime in which operational monitoring actually takes
place, and it is the regime in which the leading indicator is the only thing left
to measure.

The pattern is the precursor relation that underlies pyramid models of accident
causation [CITATION NEEDED: Heinrich's accident pyramid], and it is the rationale
for estimating rare-event frequencies from precursors \cite{bieryi1995}. The
results also show why the ratio in such models cannot be treated as a constant.
The number of violations per collision shifts by two to three orders of magnitude
between the two configurations, from 3.8 to at least 1,464 for GAP\_TTC, from 1.6
to 527 for PED\_CPA and from 22.2 to at least 2,562 for REAR\_TTC. A fixed ratio
between precursors and losses is therefore not a property that can be assumed for
a given system; it has to be measured for the configuration in question, which is
also the conclusion of the empirical critique of the fixed-ratio pyramid
\cite{yoriomoore2018}.

% =============================================================================
\section{Exposure against Configuration}
\label{sec:disc_exposure}

Every leading SPI in the register is associated with the collisions of its own
chain, yet they do not all respond to a change in the configuration. The
distinction matters because an indicator that is attached to a safety case is
expected to move when the argued property changes.

PED\_CPA responds. Its violation count falls from 1,268 in Campaign~1 to 527 in
Campaigns~2 and~3, and the pedestrian collisions fall with it, from 814 to one.
GAP\_TTC and REAR\_TTC do not. GAP\_TTC is violated 2,070 times in Campaign~1 and
5,624 times in Campaigns~2 and~3, and REAR\_TTC 2,243 and 9,841 times, while the
collisions of both chains fall from 540 and 101 to zero. Allowing for the
different eligible populations, 15,362 against 65,093 runs for the oncoming chain
and 14,373 against 75,509 for the following-vehicle chain, the counts of these
two indicators are dominated by how often the conflict is encountered rather than
by how it is handled. Both indicators identify the dangerous runs within a
configuration, as their recall of 97.6\,\% and 100\,\% shows, and both are poor
signals of a change between configurations.

The same split appears when the violation rates of the two configurations are
set against each other, which removes the difference in eligible populations.
The violation rate of PED\_CPA in Campaign~1 is 8.25 times its rate in
Campaigns~2 and~3, whereas the corresponding ratio is 1.56 for GAP\_TTC and 1.2
for REAR\_TTC. REQ\_DECEL, AEB\_ONC and AEB\_PED lie between these extremes, at
3.59, 3.54 and 4.44. An indicator whose rate barely moves while the collisions of
its chain fall from several hundred to zero cannot, on its own, signal that the
argued property has changed, however well it separates dangerous from uneventful
runs within one configuration.

This distinction carries over directly to the SPI threshold T defined in
Section~\ref{sec:method_spithreshold}. Because T bounds the violation count of a
population of runs against a baseline rate, an indicator whose rate is dominated
by exposure moves little against T when the configuration changes, whereas an
indicator whose rate tracks the configuration moves strongly. The ratios above
therefore indicate which entries of the register are suited to rate-level
monitoring and which are suited mainly to the per-run detection reported in
Section~\ref{sec:results_statistics}.

% =============================================================================
\section{Geometry Decides the Warning Time}
\label{sec:disc_geometry}

The lead times in Table~\ref{tab:lead_time} divide the register along the
conflict geometry rather than along the indicator design, which is the most
consequential finding for the timeliness question of
Section~\ref{sec:intro_objective}.

In the two crossing conflicts the partner moves across the ego vehicle's line of
sight. Time to collision and time to the closest point of approach both depend on
the closing component of the relative motion, so they change slowly while the
paths are still separated and fall quickly only once the paths almost meet. The
median lead times of 0.60\,s for GAP\_TTC and 0.50\,s for PED\_CPA follow from
that geometry, and so does the fact that only 38 of 540 oncoming-vehicle and two
of 815 pedestrian collisions are preceded by a violation at least one second
ahead. In the following conflict the partner closes along the line of sight, the
same quantity changes monotonically throughout the approach, and 87 of 101
collisions are preceded by a violation at least one second ahead. The dependence
of a criticality metric on the geometry it is applied to is the point made in the
review of such metrics \cite{westhofen2023}, and the lead times measured here are
a quantitative instance of it.

Replacing the pedestrian time to collision by the time to the closest point of
approach was a consequence of Step~4 of the derivation, which selects the
formulation the geometry admits. The change removed alarms for pedestrians whose
path misses the vehicle, which is visible in the precision of PED\_CPA, the
highest of the register at 44.0\,\%. It did not lengthen the warning time,
because it does not change what the metric depends on. Any substantially earlier
warning for a crossing pedestrian would have to come from signals of intention
such as gait, head pose or dwell at the kerb, and those are perception
quantities, outside the behavioural tier that this work evaluates.

% =============================================================================
\section{The Pedestrian Collision Rule}
\label{sec:disc_pedrule}

The rule of Section~\ref{sec:exp_collision}, under which a pedestrian contact
counts as a collision only when the ego vehicle is moving, is a decision about
the reference event rather than about an indicator, and it deserves to be
defended as such.

Without the rule the two realistic-driving campaigns would appear to produce 481
pedestrian collisions rather than one. The appearance would be an artefact: 480
of those 481 contacts occur with a standing ego vehicle, which a CARLA walker
does not avoid. The low-collision configuration would then look dangerous for a
reason that has nothing to do with the driving behaviour under observation, and
every pedestrian indicator would be scored against the wrong events. A lagging
indicator that counts the wrong events corrupts every leading indicator derived
from it, so the first step of the derivation carries at least as much weight as
the steps that follow it.

The rule is nonetheless a judgement. The audit separates the two populations
cleanly, by ego speed, by the position of the pedestrian and by a contact impulse
that differs by about a factor of ten, and the rule is stated explicitly so that
it can be challenged. The threshold of 0.3\,m/s is still a choice, and a
sensitivity analysis under the unrestricted definition would bound how far the
pedestrian results depend on it.

% =============================================================================
\section{AEB\_PED and the Order of the Reaction Tier}
\label{sec:disc_aebped}

Two results concerning the reaction tier run against the expectation under which
it was placed in the derivation, and both are discussed here rather than treated
as settled.

AEB\_PED is the weakest entry in the register once contacts with a standing ego
vehicle are excluded. It precedes 110 of the 815 pedestrian collisions, its
relative risk is 3.7, and within Campaign~1 alone it is 1.4, which is close to no
association. An emergency brake with a pedestrian in conflict therefore carries
little information about whether that run ends in a pedestrian collision. The
reading that follows from the collision rule is that the indicator previously
appeared stronger because braking brings the vehicle to a stop and a stopped
vehicle in a walker's path is walked into, so a large part of its apparent
association was with contacts that are no longer counted.

The lead time of AEB\_PED has to be read with the same reserve. In the 110 runs
in which an event precedes a collision, it does so by between 4.95\,s and
6.24\,s, with a median of 5.15\,s. Taken at face value this would be by far the
longest warning in the register. The most likely explanation is a stop-and-go
sequence: the vehicle brakes for one pedestrian, stops, drives on, and collides
with a pedestrian later in the run, with the lead time measured from the first
violation anywhere in the run. That explanation has not been verified. Confirming
it requires a per-run check of whether the ego speed falls to near zero after the
braking event and rises again before contact, and whether the pedestrian struck
is the one braked for. Until that check is carried out the interpretation stands
as the most likely reading and not as a finding, and the figure of 5.15\,s is not
reported as a warning time.

The second anomaly concerns the ordering of the tiers. The reaction tier was
placed nearest the loss, so its lead time was expected to be the shortest of the
chain. In the oncoming chain the opposite holds: the median lead of AEB\_ONC is
0.75\,s against 0.60\,s for GAP\_TTC. Two explanations are plausible and neither
is tested here. The brake command that triggers AEB\_ONC comes from the
pedestrian-stop rule of the ego configuration and is therefore not issued in
response to the oncoming vehicle at all, so it can fire early for an unrelated
reason. Alternatively the tenth-percentile metric threshold of GAP\_TTC is tight
enough to delay the kinematic violation past the braking event. Distinguishing
the two would require the same kind of per-run attribution as the AEB\_PED check.

% =============================================================================
\section{REQ\_DECEL and the Second Branch}
\label{sec:disc_reqdecel}

REQ\_DECEL precedes 317 of the 540 oncoming-vehicle collisions, markedly fewer
than the 527 of GAP\_TTC. This is the expected consequence of the decomposition
rather than a weakness of the indicator.

The two indicators were derived from the two branches of an OR gate. GAP\_TTC
watches the accepted gap, which is the branch that applies whenever the gap was
too small from the outset. REQ\_DECEL watches the deceleration that would be
required to stop short, which is the branch that applies when the closing speed
exceeds what the turn arc allows. Each branch alone is sufficient for the
collision, so neither indicator can be expected to cover the collisions of the
other.

The formulation of REQ\_DECEL carries two assumptions that make it optimistic.
The required deceleration treats the conflict partner as stationary, so a partner
that is itself closing is not accounted for, and the range entering the
denominator is clamped at a minimum value so that the expression stays finite at
short range. The clamp produces a cluster of values near 5.41\,m/s\textsuperscript{2},
which is visible in the distribution of the metric and is an artefact of the
formulation rather than a property of the conflicts. Both assumptions bias the
required deceleration downwards and therefore make violations less frequent than
a partner-aware formulation would produce.

% =============================================================================
\section{The Choice of Metric Threshold}
\label{sec:disc_threshold}

The empirical percentile was selected in Section~\ref{sec:method_thresholds} for
reasons of independence rather than of performance, and the results allow both
sides of that choice to be stated.

Its benefit is that no metric threshold is fitted on the collisions against which
the indicators are later scored. The only collision labels available are the ones
that the evaluation uses, so any criterion that optimises a detection measure on
labelled data would have produced a circular evaluation. The percentile uses the
uneventful runs of Campaign~1 only, and the resulting thresholds were applied
unchanged to all 94,816 runs, including the two campaigns that did not
contribute to the calibration.

Its cost is equally explicit and is visible in the results. A tenth-percentile
threshold is crossed by about one tenth of uneventful runs by construction, which
is the origin of the false-alarm rate of 8.97\,\% for GAP\_TTC and 13.35\,\% for
REAR\_TTC. The metric threshold marks behaviour that is unusual relative to the
calibration population; it does not mark behaviour that is unacceptable, and
nothing in the calibration establishes a level of risk. The same holds for the
SPI threshold, which is an exact binomial limit on departure from a measured
baseline and not a risk-acceptance criterion. Deriving either from ALARP, GAMAB,
MEM or a careful and competent human driver benchmark would require inputs that
this work does not have.

% =============================================================================
\section{Limitations}
\label{sec:disc_limitations}

The preceding sections qualify individual results. This section collects the
limitations that bound the evaluation as a whole;
Table~\ref{tab:limitations} states each with its consequence and the step that
would address it.

% ---------------------------------------------------------------------------
\begin{table}[htbp]
    \centering
    \caption{Limitations of the evaluation. The table states each limitation,
    the consequence it has for the results and the step that would address it.}
    \label{tab:limitations}
    \footnotesize
    \setlength{\tabcolsep}{4pt}
    \begin{tabular}{@{}>{\raggedright\arraybackslash}p{3.4cm}>{\raggedright\arraybackslash}p{5.6cm}>{\raggedright\arraybackslash}p{4.6cm}@{}}
        \toprule
        {\bfseries Limitation} & {\bfseries Consequence} & {\bfseries Next step} \\
        \midrule
        Traffic Manager ego without gap acceptance &
        Oncoming conflicts arise from a missing capability rather than from a
        functional insufficiency in the sense of ISO~21448 \cite{iso21448} &
        A gap-accepting driving function \\
        Ground-truth signals &
        The independence of the indicators from perception holds in simulation
        only &
        Perception in the loop; architectural-tier indicators \\
        Narrow exposure &
        One scenario, one map, oncoming partners that do not yield &
        Naturalistic data; responsive partners \\
        Statistical thresholds &
        $\theta$ and T bound demonstrated behaviour; ALARP, GAMAB, MEM and the
        careful and competent human driver benchmark are named but not applied &
        A risk-acceptance SPI threshold set with safety engineers \\
        Campaign comparison &
        The campaigns differ in scenario generator as well as controller, so the
        comparison is one of known groups and not a controlled experiment &
        Controlled variation of one factor \\
        Lead-time definition &
        Lead time is measured from the first violation anywhere in the run &
        Attribution of violations to the conflict partner \\
        \bottomrule
    \end{tabular}
\end{table}

The last two limitations bear on specific results. Because the campaigns differ
in two respects at once, the contrast of Section~\ref{sec:disc_exposure} shows
that PED\_CPA tracks the difference between two configurations, not that it
tracks the controller change within it. Because lead time is measured from the
first violation anywhere in the run, the AEB\_PED interval of
Section~\ref{sec:disc_aebped} cannot be attributed to the collision that follows
it without the per-run check described there.